\section{Algorithm-Level Implementation Details}
\label{section_implementation}

    This section details the algorithm-level implementation choices that keep the online optimization lightweight: execution-time prediction and hyperparameters.

    \subsection{Predicting Execution-Time Distributions}
    \label{section: predict_ET_exp}
    We assume the robot's working environment changes smoothly, so over a short prediction interval a task's execution-time distribution does not change much.
    We therefore predict each task's execution-time distribution for the upcoming interval directly from its execution times measured at the start of that interval.
    This simple heuristic trades prediction accuracy for low online computation cost, which suits our resource-constrained setting.
    More sophisticated prediction models are possible and left for future work.


    % --- Mapping Priority Assignments to Priority Values: commented out of main text, retained for future reference ---
    % \subsection{Mapping Priority Assignments to Priority Values}
    % \label{section_pa_values}
    % While the priority assignment algorithm determines the relative priority order of tasks, Linux requires specific \textit{priority values} (real-time processes range from 0 to 99, larger is higher) to implement these assignments.
    % Theoretically, any mapping that adheres to the priority assignment order is valid.
    % However, changing priority values involves numerous low-level OS operations, and this overhead grows for multi-threaded tasks whose threads must update simultaneously, so we minimize the number of value changes.
    % We select the task with the largest number of threads as the \textit{reference task} and assign it a fixed priority value; all other tasks' values are computed relative to it, following the priority assignment order.
    % This minimizes changes to the reference task's priority value; if another task's thread count surpasses the reference task's at runtime, that task is dynamically designated as the new reference.
    % The strategy generalizes to multi-core systems under partitioned scheduling, where a fixed priority value can be locked for one task per core.
    %
    % \begin{Example}
    % Suppose the scheduler determines the priority order MPC $>$ RRT $>$ TSP $>$ SLAM, and SLAM has the highest number of threads.
    % Assigning SLAM a fixed priority value of 5, an acceptable mapping for the other tasks is MPC=8, RRT=7, TSP=6, SLAM=5.
    % \end{Example}

    \subsection{Hyperparameter Configuration}
    Only a few hyperparameters require configuration.

    Distribution resolution is the number of value-probability pairs describing each distribution during the fixed-point iteration in~\eqref{eq_prob_rta}. We use 10, beyond which no observable improvement justifies the extra cost.

    Beam width is the number of partial priority assignments kept during priority assignment optimization ($m$ in Algorithm~\ref{alg:cabs}, Section~\ref{section_cabs}). We set it to 2 to keep overall overhead below 1\%.

    Time-limit search patience bounds the QoS-budget walk (Section~\ref{section_qos_walk}): it counts consecutive non-improving steps before a pass halts, so a larger value permits a wider search at higher online cost. We use 0 for incremental intervals and 1 for re-optimization intervals started from a fresh seed, which need a slightly wider search to recover.

    Scheduler period $T$ is the interval between scheduler invocations. We recommend 10~s, the lowest value keeping scheduler overhead below 1\%.
